\documentclass{article}

% if you need to pass options to natbib, use, e.g.:
%     \PassOptionsToPackage{numbers, compress}{natbib}
% before loading neurips_2025

% The authors should use one of these tracks.
% Before accepting by the NeurIPS conference, select one of the options below.
% 0. "default" for submission
 \usepackage{neurips_2025}
% the "default" option is equal to the "main" option, which is used for the Main Track with double-blind reviewing.
% 1. "main" option is used for the Main Track
%  \usepackage[main]{neurips_2025}
% 2. "position" option is used for the Position Paper Track
%  \usepackage[position]{neurips_2025}
% 3. "dandb" option is used for the Datasets & Benchmarks Track
 % \usepackage[dandb]{neurips_2025}
% 4. "creativeai" option is used for the Creative AI Track
%  \usepackage[creativeai]{neurips_2025}
% 5. "sglblindworkshop" option is used for the Workshop with single-blind reviewing
 % \usepackage[sglblindworkshop]{neurips_2025}
% 6. "dblblindworkshop" option is used for the Workshop with double-blind reviewing
%  \usepackage[dblblindworkshop]{neurips_2025}

% After being accepted, the authors should add "final" behind the track to compile a camera-ready version.
% 1. Main Track
 % \usepackage[main, final]{neurips_2025}
% 2. Position Paper Track
%  \usepackage[position, final]{neurips_2025}
% 3. Datasets & Benchmarks Track
 % \usepackage[dandb, final]{neurips_2025}
% 4. Creative AI Track
%  \usepackage[creativeai, final]{neurips_2025}
% 5. Workshop with single-blind reviewing
%  \usepackage[sglblindworkshop, final]{neurips_2025}
% 6. Workshop with double-blind reviewing
%  \usepackage[dblblindworkshop, final]{neurips_2025}
% Note. For the workshop paper template, both \title{} and \workshoptitle{} are required, with the former indicating the paper title shown in the title and the latter indicating the workshop title displayed in the footnote.
% For workshops (5., 6.), the authors should add the name of the workshop, "\workshoptitle" command is used to set the workshop title.
% \workshoptitle{WORKSHOP TITLE}

% "preprint" option is used for arXiv or other preprint submissions
 % \usepackage[preprint]{neurips_2025}

% to avoid loading the natbib package, add option nonatbib:
%    \usepackage[nonatbib]{neurips_2025}

\usepackage[utf8]{inputenc} % allow utf-8 input
\usepackage[T1]{fontenc}    % use 8-bit T1 fonts
\usepackage{hyperref}       % hyperlinks
\usepackage{url}            % simple URL typesetting
\usepackage{booktabs}       % professional-quality tables
\usepackage{amsfonts}       % blackboard math symbols
\usepackage{nicefrac}       % compact symbols for 1/2, etc.
\usepackage{microtype}      % microtypography
\usepackage{xcolor}         % colors

% Note. For the workshop paper template, both \title{} and \workshoptitle{} are required, with the former indicating the paper title shown in the title and the latter indicating the workshop title displayed in the footnote. 
\title{Formatting Instructions For NeurIPS 2025}


% The \author macro works with any number of authors. There are two commands
% used to separate the names and addresses of multiple authors: \And and \AND.
%
% Using \And between authors leaves it to LaTeX to determine where to break the
% lines. Using \AND forces a line break at that point. So, if LaTeX puts 3 of 4
% authors names on the first line, and the last on the second line, try using
% \AND instead of \And before the third author name.


\author{%
  David S.~Hippocampus\thanks{Use footnote for providing further information
    about author (webpage, alternative address)---\emph{not} for acknowledging
    funding agencies.} \\
  Department of Computer Science\\
  Cranberry-Lemon University\\
  Pittsburgh, PA 15213 \\
  \texttt{hippo@cs.cranberry-lemon.edu} \\
  % examples of more authors
  % \And
  % Coauthor \\
  % Affiliation \\
  % Address \\
  % \texttt{email} \\
  % \AND
  % Coauthor \\
  % Affiliation \\
  % Address \\
  % \texttt{email} \\
  % \And
  % Coauthor \\
  % Affiliation \\
  % Address \\
  % \texttt{email} \\
  % \And
  % Coauthor \\
  % Affiliation \\
  % Address \\
  % \texttt{email} \\
}

\usepackage{graphicx}
\usepackage{apptools}
\usepackage[flushleft]{threeparttable}
\usepackage{array,booktabs,makecell}
\usepackage{multirow}
\usepackage{nicefrac}
\usepackage{amsthm}
\usepackage{mathtools}
\usepackage{amsmath,amsfonts,bm,amssymb}
\usepackage{dsfont}
\usepackage{wrapfig}
\usepackage{caption}
% \usepackage{subcaption}
\usepackage{siunitx}
\usepackage{thm-restate}
\usepackage{nccmath}
\usepackage{empheq}
\usepackage{bbm}
\usepackage{tabularx}

% algorithmic and algorithm already loaded by icml2026.sty
\usepackage{algorithmic}
\usepackage{algorithm}
\usepackage{filecontents}

\usepackage[capitalize,noabbrev]{cleveref}

% COLORS
\usepackage{color}
\usepackage{colortbl}
\definecolor{bgcolor}{rgb}{0.76,0.88,0.50}
\definecolor{bgcolor0}{rgb}{0.93,0.99,1}
\definecolor{bgcolor1}{rgb}{0.8,1,1}
\definecolor{bgcolor2}{rgb}{0.8,1,0.8}
\definecolor{bgcolor3}{rgb}{0.50,0.90,0.50}
\usepackage{tcolorbox}
\usepackage{pifont}

\definecolor{mydarkblue}{rgb}{0,0.16,0.54}
\definecolor{mydarkgreen}{RGB}{39,130,67}
\definecolor{mydarkorange}{RGB}{236,147,14}
% \definecolor{myorange}{RGB}{230,120,20}
% \definecolor{mydarkred}{RGB}{192,47,25}
\definecolor{mydarkgreen}{RGB}{0,100,0}
\definecolor{ruby}{RGB}{155,17,30}
\definecolor{chili}{RGB}{191,0,0}
\definecolor{sangria}{RGB}{146,0,10}
\definecolor{burgundy}{RGB}{128,0,32}
\definecolor{darkred}{RGB}{132,0,0} 
\definecolor{cherry}{RGB}{192,0,0} 
\definecolor{myaccent}{rgb}{0,0.45,0.45}
\definecolor{anotherdarkred}{rgb}{0.45,0,0.05}
\newcommand{\green}{\color{mydarkgreen}}
\newcommand{\orange}{\color{mydarkorange}}
\newcommand{\blue}{\color{mydarkblue}}
\newcommand{\red}{\color{cherry}}
% \newcommand{\red}{\color{anotherdarkred}}
\newcommand{\highlightcolor}{\color{myaccent}}

\newcommand{\markchanges}{}

% todonotes already loaded in paper.tex
\usepackage[textsize=tiny]{todonotes}
\newcommand{\arto}[1]{\todo[inline]{\textbf{Arto: }#1}}
\newcommand{\zhiro}[1]{\todo[inline]{\textbf{Zhirayr: }#1}}
\newcommand{\peter}[1]{\todo[inline]{\textbf{Peter: }#1}}

\usepackage{xspace}
\usepackage{relsize}
\newcommand{\dataname}[1]{{\tt\footnotesize\color{blue}#1}\xspace}


\newcommand{\ind}{\perp\!\!\!\!\perp}
\newcommand{\norm}[1]{\left\| #1 \right\|}
\newcommand{\sqnorm}[1]{\left\| #1 \right\|^2}
\newcommand{\infnorm}[1]{\left\| #1 \right\|_{\infty}}
\newcommand{\flr}[1]{\left\lfloor #1\right\rfloor} % floor
\newcommand{\ceil}[1]{\left\lceil #1\right\rceil} % ceiling
\newcommand{\lin}[1]{\left\langle #1\right\rangle} % inner product
\newcommand{\inp}[2]{\left\langle#1,#2\right\rangle} % inner product
\newcommand{\abs}[1]{\left| #1 \right|}

\newcommand{\R}{\mathbb{R}} % reals
\newcommand{\N}{\mathbb{N}} % natural numbers
\newcommand{\Z}{\mathbb{Z}} % integers
\newcommand{\E}[1]{\mathbb{E}\left[#1\right]}
\newcommand{\Med}[1]{\mathrm{Med}\left[#1\right]}
\newcommand{\Var}[1]{\mathrm{Var}\left(#1\right)}

\newcommand{\supp}{\operatorname{supp}}
\newcommand{\progg}{\mathrm{prog}}

\newcommand{\Exp}[1]{{\mathbb{E}}\left[#1\right]}
\newcommand{\ExpSub}[2]{{\mathbb{E}}_{#1}\left[#2\right]}
\newcommand{\ExpCond}[2]{{\mathbb{E}}\left[\left.#1\ \right\vert\ #2\right]}

\newcommand{\Prob}[1]{\mathbb{P}\left(#1\right)} % probability
\newcommand{\ProbCond}[2]{\mathbb{P}\left(#1\middle\vert#2\right)}

% hat
\newcommand{\hmu}{\hat{\mu}}

% confidence
\newcommand{\conf}{\mathrm{conf}}

% caligraphic
\newcommand{\cA}{\mathcal{A}}
\newcommand{\cB}{\mathcal{B}}
\newcommand{\cC}{\mathcal{C}}
\newcommand{\cD}{\mathcal{D}}
\newcommand{\cE}{\mathcal{E}}
\newcommand{\cF}{\mathcal{F}}
\newcommand{\cG}{\mathcal{G}}
\newcommand{\cH}{\mathcal{H}}
\newcommand{\cI}{\mathcal{I}}
\newcommand{\cJ}{\mathcal{J}}
\newcommand{\cK}{\mathcal{K}}
\newcommand{\cL}{\mathcal{L}}
\newcommand{\cM}{\mathcal{M}}
\newcommand{\cN}{\mathcal{N}}
\newcommand{\cO}{\mathcal{O}}
\newcommand{\cP}{\mathcal{P}}
\newcommand{\cQ}{\mathcal{Q}}
\newcommand{\cR}{\mathcal{R}}
\newcommand{\cS}{\mathcal{S}}
\newcommand{\cT}{\mathcal{T}}
\newcommand{\cU}{\mathcal{U}}
\newcommand{\cV}{\mathcal{V}}
\newcommand{\cW}{\mathcal{W}}
\newcommand{\cX}{\mathcal{X}}
\newcommand{\cY}{\mathcal{Y}}
\newcommand{\cZ}{\mathcal{Z}}

% bold letters
\newcommand{\ba}{\bm{a}}
\newcommand{\bb}{\bm{b}}
\newcommand{\bc}{\bm{c}}
\newcommand{\bd}{\bm{d}}
\newcommand{\be}{\bm{e}}
\newcommand{\bg}{\bm{g}}
\newcommand{\bh}{\bm{h}}
\newcommand{\bi}{\bm{i}}
\newcommand{\bj}{\bm{j}}
\newcommand{\bk}{\bm{k}}
\newcommand{\bl}{\bm{l}}
\newcommand{\bn}{\bm{n}}
\newcommand{\bo}{\bm{o}}
\newcommand{\bp}{\bm{p}}
\newcommand{\bq}{\bm{q}}
\newcommand{\br}{\bm{r}}
\newcommand{\bs}{\bm{s}}
\newcommand{\bt}{\bm{t}}
\newcommand{\bu}{\bm{u}}
\newcommand{\bv}{\bm{v}}
\newcommand{\bw}{\bm{w}}
\newcommand{\bx}{\bm{x}}
\newcommand{\by}{\bm{y}}
\newcommand{\bz}{\bm{z}}

% bold matrices
\newcommand{\mA}{\mathbf{A}}
\newcommand{\mB}{\mathbf{B}}
\newcommand{\mC}{\mathbf{C}}
\newcommand{\mD}{\mathbf{D}}
\newcommand{\mE}{\mathbf{E}}
\newcommand{\mF}{\mathbf{F}}
\newcommand{\mG}{\mathbf{G}}
\newcommand{\mH}{\mathbf{H}}
\newcommand{\mI}{\mathbf{I}}
\newcommand{\mJ}{\mathbf{J}}
\newcommand{\mK}{\mathbf{K}}
\newcommand{\mL}{\mathbf{L}}
\newcommand{\mM}{\mathbf{M}}
\newcommand{\mN}{\mathbf{N}}
\newcommand{\mO}{\mathbf{O}}
\newcommand{\mP}{\mathbf{P}}
\newcommand{\mQ}{\mathbf{Q}}
\newcommand{\mR}{\mathbf{R}}
\newcommand{\mS}{\mathbf{S}}
\newcommand{\mT}{\mathbf{T}}
\newcommand{\mU}{\mathbf{U}}
\newcommand{\mV}{\mathbf{V}}
\newcommand{\mW}{\mathbf{W}}
\newcommand{\mX}{\mathbf{X}}
\newcommand{\mY}{\mathbf{Y}}
\newcommand{\mZ}{\mathbf{Z}}

\newcommand{\eqdef}{\coloneqq} 
\newcommand{\Mod}[1]{\ (\mathrm{mod}\ #1)}
\newcommand{\minimize}{\mathop{\mathrm{minimize}}}


\newcommand{\tauavg}{\tau_{\mathrm{avg}}}

\newcommand{\hnabla}{\widehat{\nabla}}


\makeatletter
\newcommand{\vast}{\bBigg@{4}}

\DeclareMathOperator*{\argmin}{arg\,min}
\DeclareMathOperator*{\argmax}{arg\,max}

\def\<{\left\langle}
\def\>{\right\rangle}
\def\[{\left[}
\def\]{\right]}
\def\({\left(}
\def\){\right)}

% The old colors
% \newcommand{\crossmarkred}{{\color{mydarkred}\ding{56}}}
% \newcommand{\checkmarkgreen}{\textbf{\color{mydarkgreen}\ding{52}}}

% % No colors
% \newcommand{\checkmarkgreen}{\textbf{\ding{52}}}
% \newcommand{\crossmarkred}{\ding{56}}

% Different colors
% \definecolor{myblue}{rgb}{0.25,0.25,0.3}
% \definecolor{myorange}{rgb}{0.25,0.25,0.3}

\definecolor{myblue}{rgb}{0.3,0.25,0.2}
\definecolor{myorange}{rgb}{0.3,0.25,0.2}
\newcommand{\checkmarkgreen}{\textbf{\color{myblue}\ding{52}}}
\newcommand{\crossmarkred}{{\color{myorange}\ding{56}}}

%%%%%%%%%%%%%%%%%%%%%%%%%%%%%%%%
% Algorithm names
%%%%%%%%%%%%%%%%%%%%%%%%%%%%%%%%
% \newcommand{\algname}[1]{{\sf\footnotesize\red#1}\xspace}
% \newcommand{\algname}[1]{{\sf\red\relscale{0.90}#1}\xspace}
\newcommand{\algname}[1]{\text{\sf\relscale{0.93}#1}\xspace}
\newcommand{\algnameS}[1]{\text{\sf\scriptsize\red#1}\xspace}
% \newcommand{\algn}{{\sf\red\relscale{0.95}Ringleader ASGD}\xspace}
\newcommand{\iasgd}{{\sf\relscale{0.95}IA\textsuperscript{2}SGD}\xspace}
\newcommand{\iasgdtitle}{IA\textsuperscript{2}SGD}
\newcommand{\ringleader}{\algname{Ringleader~ASGD}}
\newcommand{\ringleadertitle}{Ringleader~ASGD}
\newcommand{\malenia}{\algname{Malenia~SGD}}
\newcommand{\maleniatitle}{Malenia~SGD}
\newcommand{\herosgd}{\algname{Hero~SGD}}
\newcommand{\herosgdtitle}{Hero~SGD}
\newcommand{\minibatch}{\algname{Minibatch~SGD}}
\newcommand{\minibatchtitle}{Minibatch~SGD}
\newcommand{\naiveminibatch}{\algname{Naive~Minibatch~SGD}}
\newcommand{\naiveminibatchtitle}{Naive~Minibatch~SGD}
\newcommand{\asyncsgdtitle}{Asynchronous~SGD}
\newcommand{\asyncsgd}{\algname{\asyncsgdtitle}}
\newcommand{\naiveasyncsgdtitle}{Naive~Asynchronous~SGD}
\newcommand{\naiveasyncsgd}{\algname{\naiveasyncsgdtitle}}
\newcommand{\sag}{\algname{SAG}}
\newcommand{\sagtitle}{SAG}
\newcommand{\saga}{\algname{SAGA}}
\newcommand{\sagatitle}{SAGA}
\newcommand{\ringmaster}{\algname{Ringmaster~ASGD}}
\newcommand{\ringmastertitle}{Ringmaster~ASGD}
\newcommand{\rennalatitle}{Rennala~SGD}
\newcommand{\rennala}{\algname{Rennala~SGD}}
\newcommand{\sgd}{\algname{SGD}}
\newcommand{\ata}{\algname{ATA}}
\newcommand{\hogwild}{\algname{HOGWILD!}}
\newcommand{\mvrtitle}{MVR}
\newcommand{\mvr}{\algname{MVR}}
\newcommand{\storm}{\algname{STORM}}
\newcommand{\rennalamvrtitle}{{\red Rennala~MVR}}
\newcommand{\rennalamvr}{\algname{\red Rennala~MVR}}


%%%%%%%%%%%%%%%%%%%%%%%%%%%%%%%%
% THEOREMS
%%%%%%%%%%%%%%%%%%%%%%%%%%%%%%%%
\usepackage{thmtools}
\usepackage{thm-restate}
\usepackage{mdframed}

\theoremstyle{plain}
\newtheorem{theorem}{Theorem}[section]
\newtheorem{proposition}[theorem]{Proposition}
\newtheorem{lemma}[theorem]{Lemma}
\newtheorem{corollary}[theorem]{Corollary}
\theoremstyle{definition}
\newtheorem{definition}[theorem]{Definition}
\newtheorem{assumption}[theorem]{Assumption}
\theoremstyle{remark}
\newtheorem{remark}[theorem]{Remark}

% \theoremstyle{definition} % make theorems body non-italic
\theoremstyle{plain} % make theorems body italic

% Box style (used for boxed versions)
\newmdenv[
  font=\normalfont\bfseries,
  linecolor=black,
  linewidth=0.8pt,
  topline=false,
  bottomline=false,
  leftline=false,
  rightline=false,
  backgroundcolor=gray!13,
  skipabove=2pt,
  skipbelow=2pt,
  innertopmargin=-5pt,
  innerbottommargin=4pt,
  innerrightmargin=4pt,
  innerleftmargin=4pt,
]{myboxed}

% Boxed versions of theorem and lemma
\newmdtheoremenv[
  font=\normalfont\bfseries,
  linecolor=black,
  linewidth=0.8pt,
  topline=false,
  bottomline=false,
  leftline=false,
  rightline=false,
  backgroundcolor=gray!13,
  skipabove=2pt,
  skipbelow=2pt,
  innertopmargin=-5pt,
  innerbottommargin=4pt,
  innerrightmargin=4pt,
  innerleftmargin=4pt,
]{boxedtheorem}[theorem]{Theorem}

\newmdtheoremenv[
  font=\normalfont\bfseries,
  linecolor=black,
  linewidth=0.8pt,
  topline=false,
  bottomline=false,
  leftline=false,
  rightline=false,
  backgroundcolor=gray!13,
  skipabove=2pt,
  skipbelow=2pt,
  innertopmargin=-5pt,
  innerbottommargin=4pt,
  innerrightmargin=4pt,
  innerleftmargin=4pt,
]{boxedlemma}[theorem]{Lemma}

\newmdtheoremenv[
  font=\normalfont\bfseries,
  linecolor=black,
  linewidth=0.8pt,
  topline=false,
  bottomline=false,
  leftline=false,
  rightline=false,
  backgroundcolor=gray!13,
  skipabove=2pt,
  skipbelow=2pt,
  innertopmargin=-5pt,
  innerbottommargin=4pt,
  innerrightmargin=4pt,
  innerleftmargin=4pt,
]{boxeddefinition}[theorem]{Definition}

\newmdtheoremenv[
  font=\normalfont\bfseries,
  linecolor=black,
  linewidth=0.8pt,
  topline=false,
  bottomline=false,
  leftline=false,
  rightline=false,
  backgroundcolor=gray!13,
  skipabove=2pt,
  skipbelow=2pt,
  innertopmargin=-5pt,
  innerbottommargin=4pt,
  innerrightmargin=4pt,
  innerleftmargin=4pt,
]{boxedassumption}[theorem]{Assumption}

% Custom environments for restating (unboxed)
\newtheorem{innercustomthm}{Theorem}
\newenvironment{restate-theorem}[1]
  {\renewcommand\theinnercustomthm{#1}\innercustomthm}
  {\endinnercustomthm}

\newtheorem{innercustomlemma}{Lemma}
\newenvironment{restate-lemma}[1]
  {\renewcommand\theinnercustomlemma{#1}\innercustomlemma}
  {\endinnercustomlemma}

\newtheorem{innercustomproposition}{Proposition}
\newenvironment{restate-proposition}[1]
  {\renewcommand\theinnercustomproposition{#1}\innercustomproposition}
  {\endinnercustomproposition}

% Boxed restated theorem and lemma
\newenvironment{restate-boxedtheorem}[1]
  {\renewcommand\theinnercustomthm{#1}\begin{myboxed}\begin{innercustomthm}}
  {\end{innercustomthm}\end{myboxed}}

\newenvironment{restate-boxedlemma}[1]
  {\renewcommand\theinnercustomlemma{#1}\begin{myboxed}\begin{innercustomlemma}}
  {\end{innercustomlemma}\end{myboxed}}

% Sketch of Proof
\newcommand*{\sketchproofname}{Sketch of Proof}
\newenvironment{sketchproof}[1][\sketchproofname]
  {\begin{proof}[#1]}
  {\end{proof}}


\begin{document}


\maketitle


% \begin{abstract}
%     The abstract paragraph should be indented \nicefrac{1}{2}~inch (3~picas) on
%     both the left- and right-hand margins. Use 10~point type, with a vertical
%     spacing (leading) of 11~points.  The word \textbf{Abstract} must be centered,
%     bold, and in point size 12. Two line spaces precede the abstract. The abstract
%     must be limited to one paragraph.
% \end{abstract}



\section{Problem Setup}
% 
% 
We consider the nonconvex optimization problem
\begin{equation}\label{eq:problem}
  \minimize_{x \in \R^d} 
  \left\{ f(x) \coloneqq \ExpSub{\xi \sim \cD}{f(x;\xi)} \right\} ,
\end{equation}
where $f(x; \xi)$ is the loss function evaluated on a data sample $\xi$ drawn from distribution~$\cD$, and the model is parameterized by $x \in \mathbb{R}^d$ with $d$ denoting the parameter dimensionality.
% The goal is to learn a global model that minimizes the expected loss.

We consider a distributed learning setting with $n$ workers, where each worker~$i$ has access to the same data distribution~$\cD$.
This setting is common in data centers with either unified memory or uniformly partitioned data.
Following the \textit{fixed computation model} \citep{mishchenko2022asynchronous}, we formalize the heterogeneous computation times as follows:
% 
\begin{equation}\label{eq:fixed_time}
  \begin{minipage}{0.9\linewidth}
    \centering
    Each worker~$i$ requires $\tau_i$~seconds to compute one stochastic~gradient~$\nabla f(x;\xi)$.

    Without loss of generality, we assume $0 < \tau_1 \le \tau_2 \le \cdots \le \tau_n$.
  \end{minipage}
\end{equation}
% 
We assume instantaneous communication (zero latency) between workers and the server in both directions.
This modeling simplification has been the standard assumption in prior work analyzing the time complexity of distributed optimization \citep{mishchenko2022asynchronous, koloskova2022sharper, tyurin2023optimal, maranjyan2025ringmaster}.


\section{Algorithm}
% 
% 
% 
\begin{algorithm}[H]
    \caption{\ringmaster (without calculation stops)}
    \label{alg:ringmasternew}
    \begin{algorithmic}[1]
        \STATE \textbf{Input:} point $x^0 \in \R^d$, stepsize $\gamma > 0,$ delay threshold $R \in \N$
        \STATE Set $k = 0$
        \STATE Workers start computing stochastic gradients at $x^0$
        \WHILE{not terminated}
            \STATE Gradient $\nabla f\big(x^{k-\delta^k}; \xi^{k-\delta^k}_{i}\big)$ arrives from worker $i$
            \highlightcolor
            \IF{$\delta^k < R$}
                \color{black}
                \STATE Update the model:
                    $x^{k+1} = x^{k} - \gamma \nabla f\big(x^{k-\delta^k}; \xi^{k-\delta^k}_{i}\big)$
                \STATE Worker $i$ begins calculating 
                    $\nabla f\big(x^{k+1}; \xi^{k+1}_{i}\big)$
                \STATE Update the iteration number $k = k + 1$
            \ELSE
                \STATE Ignore the outdated gradient 
                    $\nabla f\big(x^{k-\delta^k}; \xi^{k-\delta^k}_{i}\big)$
                \STATE Worker $i$ begins calculating 
                    $\nabla f\big(x^{k}; \xi^{k}_{i}\big)$
            \ENDIF
        \ENDWHILE
    \end{algorithmic}
\end{algorithm}
% 
% 
\begin{algorithm}[H]
    \caption{\ringmaster (with calculation stops)}
    \label{alg:ringmasternewstops}
    \begin{algorithmic}[1]
        \STATE \textbf{Input:} point $x^0 \in \R^d$, stepsize $\gamma > 0$, delay threshold $R \in \N$
        \STATE Set $k = 0$
        \STATE Workers start computing stochastic gradients at $x^0$
        \WHILE{not terminated}
            \STATE {\highlightcolor Stop calculating stochastic gradients with delays $\geq R$, and start computing new ones at $x^k$ instead}
            \STATE Gradient $\nabla f\big(x^{k-\delta^k}; \xi^{k-\delta^k}_{i}\big)$ arrives from worker $i$
            \STATE Update the model:
                $x^{k+1} = x^{k} - \gamma \nabla f\big(x^{k-\delta^k}; \xi^{k-\delta^k}_{i}\big)$
            \STATE Worker $i$ begins calculating 
                $\nabla f\big(x^{k+1}; \xi^{k+1}_{i}\big)$
            \STATE Update the iteration number $k = k + 1$
        \ENDWHILE
    \end{algorithmic}
\end{algorithm}
% 
% 
\section{Lemmas}

\begin{lemma}[{\citet[Lemma 4.1]{maranjyan2025ringmaster}}]
    \label{lem:time_R}
    Let the workers' computation times satisfy the \emph{fixed computation model} \eqref{eq:fixed_time}.
    Let $R$ be the delay threshold of \Cref{alg:ringmasternew} or \Cref{alg:ringmasternewstops}.
    The time required to complete any $R$ consecutive iterate updates of \Cref{alg:ringmasternew} or \Cref{alg:ringmasternewstops} is at most
    \begin{equation}\label{eq:time_R}
       t(R) 
       \eqdef 2 \min\limits_{m \in [n]} \left(\frac{1}{m} \sum\limits_{i=1}^m \frac{1}{\tau_i}\right)^{-1} \left( 1 + \frac{R}{m} \right).
    \end{equation}
\end{lemma}

So the time complexity of the algorithm with fixed threshold $R$ to do $K$ iterations is
$$
    \frac{K}{R}\times t(R).
$$
% 
\section{Time Complexity for Changing Threshold}
% 
In this section, we provide a formal derivation for the time complexity of the algorithm when the delay threshold $R_k$ is a non-decreasing function of the iteration index $k$.
% 
\begin{lemma}[Time Complexity for Adaptive Threshold]
    \label{lem:time_RK}
    Let the workers' computation times satisfy the \textit{fixed computation model} \eqref{eq:fixed_time}. 
    Let $R_k$ be a non-decreasing delay threshold for $k \in \{0, 1, \dots, K-1\}$.
    For each integer value $r \in \{\lfloor R_0 \rfloor, \dots, \lfloor R_{K-1} \rfloor\}$, define the set of iterations $\mathcal{K}_r = \{k \mid \lfloor R_k \rfloor = r\}$.
    The total time $T(K)$ required to complete $K$ iterations is bounded by
    \begin{equation*}
        T(K) \le \sum_{r=\lfloor R_0 \rfloor}^{\lfloor R_{K-1} \rfloor} \left\lceil \frac{|\mathcal{K}_r|}{r} \right\rceil t(r) ,
    \end{equation*}
    where $t(r) = 2 \min_{m \in [n]} H_m \left( 1 + \frac{r}{m} \right)$ as defined in Lemma~\ref{lem:time_R}, and $H_m = \left(\frac{1}{m} \sum_{i=1}^m \frac{1}{\tau_i}\right)^{-1}$.
\end{lemma}
% 
% 
\begin{proof}
    We partition the set of iterations $\{0, 1, \dots, K-1\}$ into disjoint blocks $\mathcal{K}_r$ based on the integer value of the threshold $R_k$. Since $R_k$ is non-decreasing, these blocks consist of consecutive iteration indices. For any specific block $\mathcal{K}_r$, the delay threshold satisfies $r\le R_k < r+1$. 
    
    By Lemma~\ref{lem:time_R}, the time required to complete any $r$ consecutive updates when the threshold is $r$ is at most $t(r)$.
    We can cover the $|\mathcal{K}_r|$ iterations in block $r$ by dividing them into $N_b$ sub-blocks of size exactly $r$, where:
    \begin{equation*}
        N_b = \left\lceil \frac{|\mathcal{K}_r|}{r} \right\rceil .
    \end{equation*}
    The last sub-block may contain fewer than $r$ updates, but its duration is still upper-bounded by $t(r)$, as the time to complete $r' < r$ updates is strictly less than the time to complete $r$ updates. Therefore, the time spent in each block $\mathcal{K}_r$ is bounded by $N_b \cdot t(r)$. Summing over all possible values of $r$ that the threshold takes during the $K$ iterations gives:
    \begin{equation*}
        T(K) = \sum_{r=\lfloor R_0 \rfloor}^{\lfloor R_{K-1} \rfloor} T(\mathcal{K}_r) \le \sum_{r=\lfloor R_0 \rfloor}^{\lfloor R_{K-1} \rfloor} \left\lceil \frac{|\mathcal{K}_r|}{r} \right\rceil t(r).
    \end{equation*}
    This completes the proof.
\end{proof}
% 
% 
\begin{corollary}[Time Complexity for Square-Root Threshold]
    \label{cor:time_sqrt}
    Let the delay threshold grow according to $R_k = \sqrt{k+1}$ for $k \ge 0$. Under the assumptions of Lemma~\ref{lem:time_RK}, the total time to reach $K$ iterations is 
    \begin{equation*}
        T(K) = \mathcal{O}\( \min_{m \in [n]}  H_m \( \sqrt{K} + \frac{K}{m} \) \).
    \end{equation*}
\end{corollary}
% 
\begin{proof}
    We first determine the size of the blocks $|\mathcal{K}_r|$. For a given integer $r$, the condition $r \le \sqrt{k+1} < r+1$ implies $r^2 \le k+1 < (r+1)^2$. The range of iteration indices $k$ belonging to $\mathcal{K}_r$ is thus:
    \begin{equation*}
        k \in \{r^2-1, r^2, \dots, (r+1)^2-2\}.
    \end{equation*}
    The number of iterations in block $r$ is:
    \begin{equation*}
        |\mathcal{K}_r| = \left( (r+1)^2 - 2 \right) - (r^2 - 1) + 1 = (r+1)^2 - r^2 = 2r + 1.
    \end{equation*}
    Next, we substitute $|\mathcal{K}_r|$ into the summation from Lemma~\ref{lem:time_RK}. For $r \ge 1$, the ceiling term is:
    \begin{equation*}
        \left\lceil \frac{2r+1}{r} \right\rceil = \left\lceil 2 + \frac{1}{r} \right\rceil = 3.
    \end{equation*}
    The total time to reach $K$ iterations, where the maximum threshold reached is $r_{\max} = \lfloor \sqrt{K} \rfloor$, is:
    \begin{equation*}
        T(K) \le \sum_{r=1}^{\lfloor \sqrt{K} \rfloor} 3 \cdot t(r) = 6 \sum_{r=1}^{\lfloor \sqrt{K} \rfloor} \min_{m \in [n]} \left[ H_m \left( 1 + \frac{r}{m} \right) \right].
    \end{equation*}
    Using the property that the minimum of a sum is greater than or equal to the sum of the minimums:
    \begin{align*}
        T(K) &\le 6 \min_{m \in [n]} \left[ \sum_{r=1}^{\lfloor \sqrt{K} \rfloor} H_m + \sum_{r=1}^{\lfloor \sqrt{K} \rfloor} \frac{H_m}{m} r \right] \\
             &= 6 \min_{m \in [n]} \left[ H_m \lfloor \sqrt{K} \rfloor + \frac{H_m}{m} \frac{\lfloor \sqrt{K} \rfloor(\lfloor \sqrt{K} \rfloor+1)}{2} \right] \\
             &\le 6 \min_{m \in [n]} \left[ H_m \sqrt{K} + \frac{H_m}{m} K \right].
    \end{align*}
    Thus,
    \begin{equation*}
        T(K) = \mathcal{O}\( \min_{m \in [n]}  H_m \( \sqrt{K} + \frac{K}{m} \) \).
    \end{equation*}
\end{proof}
%%%%%%%%%%%%%%%%%%%%%%%%%%%%%%%%%%%%%%%%%%%%%%%%%%%%%%%%%%%%%%%%%%%%%%%%%%%%%%%%%%%%%%%%%%%%%%%%%%%%%%%%%%%%%%%%%%%%%%%%
%%%%%%%%%%%%%%%%%%%%%%%%%%%%%%%%%%%%%%%%%%%%%%%%%%%%%%%%%%%%%%%%%%%%%%%%%%%%%%%%%%%%%%%%%%%%%%%%%%%%%%%%%%%%%%%%%%%%%%%%
%%%%%%%%%%%%%%%%%%%%%%%%%%%%%%%%%%%%%%%%%%%%%%%%%%%%%%%%%%%%%%%%%%%%%%%%%%%%%%%%%%%%%%%%%%%%%%%%%%%%%%%%%%%%%%%%%%%%%%%%
{
\small
\bibliographystyle{apalike}
\bibliography{bib}
}
\newpage
\appendix
%%%%%%%%%%%%%%%%%%%%%%%%%%%%%%%%%%%%%%%%%%%%%%%%%%%%%%%%%%%%%%%%%%%%%%%%%%%%%%%%%%%%%%%%%%%%%%%%%%%%%%%%%%%%%%%%%%%%%%%%
%%%%%%%%%%%%%%%%%%%%%%%%%%%%%%%%%%%%%%%%%%%%%%%%%%%%%%%%%%%%%%%%%%%%%%%%%%%%%%%%%%%%%%%%%%%%%%%%%%%%%%%%%%%%%%%%%%%%%%%%
%%%%%%%%%%%%%%%%%%%%%%%%%%%%%%%%%%%%%%%%%%%%%%%%%%%%%%%%%%%%%%%%%%%%%%%%%%%%%%%%%%%%%%%%%%%%%%%%%%%%%%%%%%%%%%%%%%%%%%%%
\section{Convergence Analysis}

We introduce the following notation: linear minimization oracle $$\texttt{lmo}(y) \eqdef \argmin_{x~:~\|x\| \le 1}\left\{\inp{y}{x} \right\}.$$

\begin{definition}
    \label{def:gen_smooth}
    Let $f:\R^d\rightarrow\R$ be differentiable. Then $f$ is symmetric $(L_0, L_1)$-smooth if there exist $L_0,L_1 \ge 0$ such that for any $x, y \in \R^d$
    \begin{equation}
        \label{eq:gen_smooth}
        \norm{\nabla f(x) - \nabla f(y)}_* \le \sup_{\theta \in [0,1]}\left(L_0 + L_1 \norm{\nabla f(\theta x +(1-\theta)y)}_*\right)\norm{x-y}.
    \end{equation}
\end{definition}

\begin{lemma} 
    \label{lem:gen_smooth}
    Let $f$ be $(L_0, L_1)$-smooth. Then for any $x,y\in \R^d$ it holds 
    \begin{equation}
        f(x) \le f(y) + \inp{\nabla f(y)}{x-y} +\frac{1}{2}\left(L_0 + L_1\norm{\nabla f(y)}_*\right)\exp\left(L_1\norm{x-y}\right)\norm{x-y}^2
    \end{equation}
\end{lemma}
\begin{algorithm}[H]
    \caption{\ringmaster (without calculation stops)}
    \label{alg:ringmaster_lmo}
    \begin{algorithmic}[1]
        \STATE \textbf{Input:} points $x^0, m^0 \in \R^d$, stepsize $\eta_k > 0,$ delay threshold $R_k \in \N$, momentum parameter $\alpha_k >0$.
        \STATE Set $k = 0$
        \STATE Workers start computing stochastic gradients at $x^0$
        \WHILE{not terminated}
            \STATE Gradient $g^k = \nabla f\big(x^{k-\delta^k}; \xi^{k-\delta^k}_{i}\big)$ arrives from worker $i$
            \highlightcolor
            \IF{$\delta^k < R_k$}
                \color{black}
                \STATE Update the mometum: $m^{k+1} = (1-\alpha_k)m^{k} + \alpha_k g^k$
                \STATE Update the model:
                    $x^{k+1} = x^{k} + \eta_k\texttt{lmo}(m^{k+1})$
                \STATE Worker $i$ begins calculating 
                    $\nabla f\big(x^{k+1}; \xi^{k+1}_{i}\big)$
                \STATE Update the iteration number $k = k + 1$
            \ELSE
                \STATE Ignore the outdated gradient 
                    $\nabla f\big(x^{k-\delta^k}; \xi^{k-\delta^k}_{i}\big)$
                \STATE Worker $i$ begins calculating 
                    $\nabla f\big(x^{k}; \xi^{k}_{i}\big)$
            \ENDIF
        \ENDWHILE
    \end{algorithmic}
\end{algorithm}

\begin{lemma}
    \label{lem:descent_lemma}
    Let $f$ be $(L_0, L_1)$-smooth. Then iterates of Algorithm~\ref{alg:ringmaster_lmo} satisfy 
    \begin{equation}
        \sum_{k=0}^{K-1}\eta_k\phi_k\norm{\nabla f(x^k)}_* \leq \Delta_0 + 2\sum^{K-1}_{k=0}\eta_k\norm{\hat{e}_{k+1}}_* +\frac{L_0}{2}\sum^{K-1}_{k=0}\exp\left(L_1\eta_k\right)\eta_k^2,
    \end{equation}
    where we define error term as $\hat{e}_{k+1} \eqdef m^{k+1} - \nabla f(x^k)$ and constatnt $\phi_k  = 1-\frac{L_1\eta_k}{2}\exp(L_1\eta_k)$. 
\end{lemma}
\begin{proof}
    By Lemma~\ref{lem:gen_smooth}, we have 
    \begin{eqnarray*}
        f(x^{k+1}) &\le& f(x^k) + \inp{\nabla f(x^k)}{x^{k+1}-x^k}\\
        && +\frac{1}{2}\left(L_0 + L_1\norm{\nabla f(x^k)}_*\right)\exp\left(L_1\norm{x^{k+1}-x^k}\right)\norm{x^{k+1}-x^k}^2\\
        &\le& f(x^k) + \eta_k\inp{\nabla f(x^k)}{\texttt{lmo}(m^{k+1})} +\frac{1}{2}\left(L_0 + L_1\norm{\nabla f(x^k)}_*\right)\exp\left(L_1\eta_k\right)\eta_k^2\\
        &=& f(x^k) + \eta_k\inp{\nabla f(x^k) - m^{k+1}}{\texttt{lmo}(m^{k+1})} + \eta_k\inp{m^{k+1}}{\texttt{lmo}(m^{k+1})} \\
        && +\frac{1}{2}\left(L_0 + L_1\norm{\nabla f(x^k)}_*\right)\exp\left(L_1\eta_k\right)\eta_k^2\\
        &\le& f(x^k) + \eta_k\norm{\nabla f(x^k) - m^{k+1}}_* - \eta_k\norm{m^{k+1}}_*  +\frac{1}{2}\left(L_0 + L_1\norm{\nabla f(x^k)}_*\right)\exp\left(L_1\eta_k\right)\eta_k^2\\
        &\le& f(x^k) + 2\eta_k\norm{\nabla f(x^k) - m^{k+1}}_* - \eta_k\norm{\nabla f(x^k)}_*  +\frac{1}{2}\left(L_0 + L_1\norm{\nabla f(x^k)}_*\right)\exp\left(L_1\eta_k\right)\eta_k^2.
    \end{eqnarray*}
\end{proof}


\begin{lemma}
    \label{lem:error_bound}
    Let $f$ be symmetric $(L_0, L_1)$-smooth. If stepsize $\eta_k$ is non-increasing and momentum parameter $\alpha_k$ defined as $\alpha_0 =1$ and $\alpha_k = \frac{1}{\sqrt{k}}$ for $k \ge 1$, then  the iterates of Algorithm~\ref{alg:ringmaster_lmo} satisfy 
    \begin{eqnarray*}
        \Exp{\norm{\hat{e}_{k+1}}_*} &\le&  \sum_{t=1}^k \prod_{j=t}^{k}(1-\alpha_j) \left(L_0 + L_1\Exp{\norm{\nabla f(x^t)}_*}\right)\exp(L_1\eta_{t-1})\eta_{t-1}\\
        && + \sum_{t=1}^k \prod_{j=t+1}^{k}(1-\alpha_j) \alpha_t (L_0 + L_1 \Exp{\norm{\nabla f(x^t)}_*})\exp\left(L_1 R_t\eta_{t-1}\right)R_t\eta_{t-1}\\
        && + {\color{red} \rho} \sigma\sqrt{\sum_{t=1}^k \prod_{j=t+1}^{k}(1-\alpha_j)^2 \alpha_t^2}.
    \end{eqnarray*}
\end{lemma}
\begin{proof}
    Using momentum update, we derive
    \begin{eqnarray*}
        \hat{e}_{k+1} &=& m^{k+1} - \nabla f(x^k)\\
        &=& (1-\alpha_k)m^k +\alpha_k g^k - \nabla f(x^k)\\
        &=& (1-\alpha_k)(m^k- \nabla f(x^{k-1})) + (1-\alpha_k)(\nabla f(x^{k-1}) - \nabla f(x^k)) +\alpha_k g^k - \alpha_k\nabla f(x^k)\\
        &=& (1-\alpha_k)(m^k- \nabla f(x^{k-1})) + (1-\alpha_k)(\nabla f(x^{k-1}) - \nabla f(x^k))\\
        && +\alpha_k (\nabla f(x^{k-\delta_k}) - \nabla f(x^k))  +\alpha_k(g^k - \nabla f(x^{k-\delta_k}).
    \end{eqnarray*}

Denoting $D_k \eqdef   \nabla f(x^{k-1}) - \nabla f(x^k)$, $S_k \eqdef \nabla f(x^{k-\delta_k}) - \nabla f(x^k)$, and $e_k = g^k - \nabla f(x^{k-\delta_k})$, we obtain
\begin{eqnarray*}
    \hat{e}_{k+1} &=& (1-\alpha_k) \hat{e}_{k} + (1-\alpha_k) D_k + \alpha_k S_k +\alpha_k e_k \\
    &=&\prod_{t = 0}^{k}(1-\alpha_t) \hat{e}_{0} + \sum_{t=0}^k \prod_{j=t+1}^{k}(1-\alpha_j)(1-\alpha_t) D_t \\
    && + \sum_{t=0}^k \prod_{j=t+1}^{k}(1-\alpha_j) \alpha_t S_t + \sum_{t=0}^k \prod_{j=t+1}^{k}(1-\alpha_j) \alpha_t e_t.
\end{eqnarray*}
Setting $\alpha_k$ as $\alpha_0 = \alpha_1 =1$ and $\alpha_k = \frac{1}{\sqrt{k}}$ for $k \ge 1$, we have 
\begin{eqnarray*}
    \hat{e}_{k+1} &=& \sum_{t=0}^k \prod_{j=t+1}^{k}(1-\alpha_j)(1-\alpha_t) D_t  + \sum_{t=0}^k \prod_{j=t+1}^{k}(1-\alpha_j) \alpha_t S_t + \sum_{t=0}^k \prod_{j=t+1}^{k}(1-\alpha_j) \alpha_t e_t\\
    &=& \underbrace{\sum_{t=1}^k \prod_{j=t+1}^{k}(1-\alpha_j)(1-\alpha_t) D_t}_{\eqdef\cT_1} + \underbrace{\sum_{t=1}^k \prod_{j=t+1}^{k}(1-\alpha_j) \alpha_t S_t}_{\eqdef \cT_2} + \underbrace{\sum_{t=1}^k \prod_{j=t+1}^{k}(1-\alpha_j) \alpha_t e_t}_{\eqdef \cT_3}.
\end{eqnarray*}

To continue the proof, we need to bound the norm of terms $\cT_1$, $\cT_2$ and $\cT_3$. By triangle inequality, we get 
\begin{eqnarray*}
    \norm{\cT_1}_* &\le& \sum_{t=1}^k \prod_{j=t+1}^{k}(1-\alpha_j)(1-\alpha_t) \norm{D_t}_*\\
    &=& \sum_{t=1}^k \prod_{j=t}^{k}(1-\alpha_j) \norm{ \nabla f(x^{t-1}) - \nabla f(x^t)}_*\\
    &\le& \sum_{t=1}^k \prod_{j=t}^{k}(1-\alpha_j) \left(L_0 + L_1\norm{\nabla f(x^t)}_*\right)\exp(L_1\eta_{t-1})\eta_{t-1}.
\end{eqnarray*}

Bounding $\norm{\cT_2}_*$, we obtain 
\begin{eqnarray*}
    \norm{\cT_2}_* &\le& \sum_{t=1}^k \prod_{j=t+1}^{k}(1-\alpha_j) \alpha_t \norm{S_t}_*\\
    &=& \sum_{t=1}^k \prod_{j=t+1}^{k}(1-\alpha_j) \alpha_t \norm{\nabla f(x^{t-\delta_t}) - \nabla f(x^t)}_*\\
    &\le& \sum_{t=1}^k \prod_{j=t+1}^{k}(1-\alpha_j) \alpha_t (L_0 + L_1 \norm{\nabla f(x^t)}_*)\exp\left(L_1 \norm{x^t - x^{t-\delta_t}}\right)\norm{x^t - x^{t-\delta_t}}\\
    &\le& \sum_{t=1}^k \prod_{j=t+1}^{k}(1-\alpha_j) \alpha_t (L_0 + L_1 \norm{\nabla f(x^t)}_*)\exp\left(L_1 \norm{\sum^{t-1}_{\tau = t-\delta_t}\eta_{\tau}\texttt{lmo}(m^{\tau+1})}\right)\norm{\sum^{t-1}_{\tau = t-\delta_t}\eta_{\tau}\texttt{lmo}(m^{\tau+1})}.
\end{eqnarray*}
 
Assuming that stepsize $\eta_k$ is non-increasing, we have 
\begin{eqnarray*}
    \norm{\cT_2}_* &\le& \sum_{t=1}^k \prod_{j=t+1}^{k}(1-\alpha_j) \alpha_t (L_0 + L_1 \norm{\nabla f(x^t)}_*)\exp\left(L_1 R_t\eta_{t-1}\right)R_t\eta_{t-1}.
\end{eqnarray*}

Bounding $\norm{\cT_3}_*$, we obtain 
\begin{eqnarray*}
    \norm{\cT_3}_* &\le& {\color{red} \rho} \norm{\sum_{t=1}^k \prod_{j=t+1}^{k}(1-\alpha_j) \alpha_t e_t}_2.
\end{eqnarray*}

Taking expectation, we have 
\begin{eqnarray*}
    \Exp{\norm{\cT_3}_*} &\le& {\color{red} \rho} \Exp{\norm{\sum_{t=1}^k \prod_{j=t+1}^{k}(1-\alpha_j) \alpha_t e_t}_2}\le {\color{red} \rho} \sqrt{\Exp{\norm{\sum_{t=1}^k \prod_{j=t+1}^{k}(1-\alpha_j) \alpha_t e_t}_2^2}}\\
    &\le& {\color{red} \rho} \sqrt{\sum_{t=1}^k \prod_{j=t+1}^{k}(1-\alpha_j)^2 \alpha_t^2 \Exp{\norm{e_t}_2^2}}\le{\color{red} \rho} \sigma\sqrt{\sum_{t=1}^k \prod_{j=t+1}^{k}(1-\alpha_j)^2 \alpha_t^2} 
\end{eqnarray*}

Thus, we have 
\begin{eqnarray*}
    \Exp{\norm{\hat{e}_{k+1}}_*} &\le&  \sum_{t=1}^k \prod_{j=t}^{k}(1-\alpha_j) \left(L_0 + L_1\Exp{\norm{\nabla f(x^t)}_*}\right)\exp(L_1\eta_{t-1})\eta_{t-1}\\
    && + \sum_{t=1}^k \prod_{j=t+1}^{k}(1-\alpha_j) \alpha_t (L_0 + L_1 \Exp{\norm{\nabla f(x^t)}_*})\exp\left(L_1 R_t\eta_{t-1}\right)R_t\eta_{t-1}\\
    && + {\color{red} \rho} \sigma\sqrt{\sum_{t=1}^k \prod_{j=t+1}^{k}(1-\alpha_j)^2 \alpha_t^2}.
\end{eqnarray*}


\end{proof}


\paragraph{Standard Smoothness Case} We set $L_1 = 0$, which implies standard smoothness of $f$. Lemma~\ref{lem:descent_lemma} implies 
\begin{eqnarray*}
    \sum_{k=0}^{K-1} \eta_k \Exp{\norm{\nabla f(x^k)}_*} &\le& \Delta_0 + 2\sum^{K-1}_{k=0}\eta_k\Exp{\norm{\hat{e}_{k+1}}_*} +\frac{L_0}{2}\sum^{K-1}_{k=0}\eta_k^2
\end{eqnarray*}

Then, by Lemma~\ref{lem:error_bound}, 
\begin{eqnarray*}
    \sum_{k=0}^{K-1} \eta_k \Exp{\norm{\nabla f(x^k)}_*} &\le& \Delta_0  +\frac{L_0}{2}\sum^{K-1}_{k=0}\eta_k^2 + 2\sum^{K-1}_{k=0}\eta_k\sum_{t=1}^k \prod_{j=t}^{k}(1-\alpha_j) L_0\eta_{t-1}\\
    && + 2\sum^{K-1}_{k=0}\eta_k\sum_{t=1}^k \prod_{j=t+1}^{k}(1-\alpha_j) \alpha_t L_0 R_t\eta_{t-1}\\
    && + 2{\color{red} \rho} \sigma\sum^{K-1}_{k=0}\eta_k\sqrt{\sum_{t=1}^k \prod_{j=t+1}^{k}(1-\alpha_j)^2 \alpha_t^2}
\end{eqnarray*}

Taking $R_t = \left\lfloor\frac{1}{\alpha_t}\right\rfloor$ and $\eta_k = \frac{\eta}{(k+1)^{\nicefrac{3}{4}}}$, we have 

\begin{eqnarray*}
    \sum_{t=1}^k \prod_{j=t+1}^{k}(1-\alpha_j)^2 \alpha_t^2 &\le& \sum_{t=1}^k \prod_{j=t+1}^{k}(1-\alpha_j) \alpha_t^2 =  \sum_{t=1}^k \prod_{j=t+1}^{k}(1-j^{-\nicefrac{1}{2}}) t^{-1}\\
    &\le& 2\exp\left(\frac{1}{1-\nicefrac{1}{2}}\right) \frac{1}{\sqrt{k+1}} \le \frac{18}{(k+1)^{\nicefrac{1}{2}}};\\
    \sum_{t=1}^k \prod_{j=t+1}^{k}(1-\alpha_j) \alpha_t  R_t\eta_{t-1} &=& \sum_{t=1}^k \prod_{j=t+1}^{k}(1-\alpha_j) \alpha_t \left\lfloor\frac{1}{\alpha_t}\right\rfloor\eta_{t-1} = \sum_{t=1}^k \prod_{j=t+1}^{k}(1-j^{-\nicefrac{1}{2}}) \frac{\left\lfloor\sqrt{t}\right\rfloor}{\sqrt{t}} \eta t^{-\nicefrac{3}{4}}\\
    &\le& \eta \sum_{t=1}^k \prod_{j=t+1}^{k}(1-j^{-\nicefrac{1}{2}})  t^{-\nicefrac{3}{4}} \le 2\eta\exp\left(\frac{1}{1-\nicefrac{1}{2}}\right) (k+1)^{\nicefrac{1}{2} - \nicefrac{3}{4}}\\
    &\le& \frac{18\eta}{(k+1)^{\nicefrac{1}{4}}};\\
    \sum_{t=1}^k \prod_{j=t}^{k}(1-\alpha_j)\eta_{t-1} &=& \eta \sum_{t=1}^k \prod_{j=t+1}^{k}(1-j^{-\nicefrac{1}{2}})  t^{-\nicefrac{3}{4}}
    \le  \frac{18\eta}{(k+1)^{\nicefrac{1}{4}}}.
\end{eqnarray*}

Thus, we obtain 
\begin{eqnarray*}
    \sum_{k=0}^{K-1} \eta_k \Exp{\norm{\nabla f(x^k)}_*} &\le& \Delta_0  +\frac{L_0}{2}\sum^{K-1}_{k=0}\eta_k^2 + 36L_0\eta\sum^{K-1}_{k=0}\frac{\eta_k}{(k+1)^{\nicefrac{1}{4}}}\\
    && + 36L_0\eta\sum^{K-1}_{k=0}\frac{\eta_k}{(k+1)^{\nicefrac{1}{4}}}\\
    && + 2{\color{red} \rho} \sigma\sum^{K-1}_{k=0}\eta_k\sqrt{\frac{18}{(k+1)^{\nicefrac{1}{2}}}}\\
    &\le& \Delta_0  +\frac{L_0\eta^2}{2}\sum^{K-1}_{k=0}\frac{1}{(k+1)^{\nicefrac{3}{2}}} + 72 L_0\eta^2 \sum^{K-1}_{k=0}\frac{1}{k+1} + 9 {\color{red} \rho} \sigma \eta \sum^{K-1}_{k=0}\frac{1}{k+1}
\end{eqnarray*}

Bounding sums as follows 
\begin{eqnarray*}
    \sum^{K-1}_{k = 0} \frac{1}{(k+1)^q} &=& \sum^{K}_{k = 1}  \frac{1}{k^q} \le \begin{cases}
        1 + \int\limits^{K}_1 \frac{dt}{t} = 1+  \log K, & \text{if } q =1\\ 
        1 + \int\limits^{K}_1 \frac{dt}{t^q} = 1+  \frac{1}{q-1} - \frac{1}{K^{q-1}} \leq \frac{q}{q-1}, & \text{if } q \geq 1\
    \end{cases},
\end{eqnarray*}
we have 
\begin{eqnarray*}
    \eta K^{\nicefrac{1}{4}}\min_{k \in \{0, \dots, K\}}\Exp{\norm{\nabla f(x^k)}_*} &\le& K\eta_K \min_{k \in \{0, \dots, K\}}\Exp{\norm{\nabla f(x^k)}_*} \le \sum_{k=0}^{K-1} \eta_k \Exp{\norm{\nabla f(x^k)}_*} \\
    &\le& \Delta_0   + 75 L_0\eta^2 (1+\log K) + 9 {\color{red} \rho} \sigma \eta (1+\log K) ,
\end{eqnarray*}
where $K\ge 2$.

Thus, we have 
\begin{eqnarray*}
    \min_{k \in \{0, \dots, K\}}\Exp{\norm{\nabla f(x^k)}_*} 
    &\le& \frac{\nicefrac{\Delta_0 }{\eta}  + 75 L_0\eta (1+\log K) + 9 {\color{red} \rho} \sigma  (1+\log K)}{ K^{\nicefrac{1}{4}}},
\end{eqnarray*}
 and iteration complexity is defined as 
 \begin{equation*}
    K = \cO\left(\frac{\left(\nicefrac{\Delta_0 }{\eta} + {\color{red} \rho} \sigma +  L_0\eta\right)^4 }{\varepsilon^4}\log \frac{1}{\varepsilon}\right).
 \end{equation*}

 To compute time complexity, we define $R_k$ as
 \begin{equation*}
    R_k = \begin{cases}
        1, &\text{if } k=0;\\
        \lfloor\sqrt{k}\rfloor, &\text{if } k\ge 1;
    \end{cases}
 \end{equation*}

Thus, by Corollary~\ref{cor:time_sqrt}, time complexity is defined as
\begin{eqnarray*}
    T(K) &=& \cO\( \min_{m \in [n]}  \left(\frac{1}{m} \sum_{i=1}^m \frac{1}{\tau_i}\right)^{-1} \( \sqrt{K} + \frac{K}{m} \) \)\\
    &=& \widetilde\cO\( \min_{m \in [n]}  \left(\frac{1}{m} \sum_{i=1}^m \frac{1}{\tau_i}\right)^{-1} \( \frac{\left(\nicefrac{\Delta_0 }{\eta} + {\color{red} \rho} \sigma +  L_0\eta\right)^2 }{\varepsilon^2}+ \frac{\left(\nicefrac{\Delta_0 }{\eta} + {\color{red} \rho} \sigma +  L_0\eta\right)^4 }{m\varepsilon^4} \) \)
\end{eqnarray*}







\end{document}